% Source: https://artofproblemsolving.com/community/c3h2278240_pixelated_circles
% Stable identity: aops-topic-2278240-q01
A \emph{pixel} is a unit square centred at a lattice point. For a positive integer $r$, the \emph{pixelated circle} of radius $r$ consists of all pixels whose centres are at distance at most $r$ from the origin; its area is the number $N(r)$ of these pixels.
\begin{enumerate}
\item[(a)] Someone lists the pixelated circle of radius $3$ as the $25$ pixels with centres $(\pm3,0)$, $(0,\pm3)$, $(0,\pm2)$, $(\pm2,0)$, $(\pm2,\pm1)$, $(\pm1,\pm2)$, $(\pm1,\pm1)$, $(\pm1,0)$, $(0,\pm1)$, $(0,0)$. Which pixels are missing, and what is $N(3)$?
\item[(b)] Show that
\[N(r)=\sum_{x=-r}^{r}\left(2\left\lfloor\sqrt{r^2-x^2}\right\rfloor+1\right).\]
\item[(c)] Compute $N(5)$.
\end{enumerate}
